\documentclass[aspectratio=169,11pt]{beamer}

\usetheme{Madrid}
\usecolortheme{default}
\usefonttheme{professionalfonts}
\setbeamertemplate{navigation symbols}{}
\setbeamertemplate{footline}[frame number]

\usepackage{amsmath, amssymb, booktabs}

\title{Labor Supply, Care Work, Fertility, And Household Allocation}
\subtitle{Gender Study --- Week 2}
\author{Teaching deck draft}
\date{}

\begin{document}

\begin{frame}
  \titlepage
\end{frame}

\begin{frame}{Central question}
\begin{itemize}
    \item How do care responsibilities, fertility, and household allocation shape labor supply and earnings trajectories?
    \item This is labor economics because care work changes employment, hours, earnings, occupations, firms, schedules, and job quality.
    \item Week 2 discipline: match each empirical object to the margin it identifies.
\end{itemize}
\end{frame}

\begin{frame}{Week 1 to Week 2 bridge}
\begin{itemize}
    \item Week 1: name the labor object before naming the mechanism.
    \item Week 2: move from person-level gaps to household and lifecycle allocation.
    \item A post-birth earnings gap can reflect care time, hours, employment, occupation, firm attachment, bargaining, policy, or norms.
\end{itemize}
\end{frame}

\begin{frame}{Labor supply inside the household}
\small
\begin{tabular}{p{0.28\linewidth} p{0.58\linewidth}}
\toprule
Object & Labor-market margin \\
\midrule
Market work & participation, hours, earnings, wage growth \\
Unpaid care & time constraints, availability for work, fatigue, schedule limits \\
Household production & substitution between parental time and market inputs \\
Fertility timing & interruptions, experience, promotion, firm continuity \\
Control over resources & labor supply, spending, care allocation, outside options \\
\bottomrule
\end{tabular}
\end{frame}

\begin{frame}{Time allocation object}
\[
T_i = h_i + a_i + \ell_i
\]
\begin{itemize}
    \item \(h_i\): market work.
    \item \(a_i\): unpaid care or home production.
    \item \(\ell_i\): leisure.
    \item A birth shock raises care demand and changes the feasible allocation of time.
\end{itemize}
\end{frame}

\begin{frame}{Household production}
\[
\max U(c_f,c_m,\ell_f,\ell_m,q)
\quad \text{s.t.} \quad
c_f+c_m+p_xx=w_fh_f+w_mh_m+y,
\quad q=F(a_f,a_m,x)
\]
\begin{itemize}
    \item Care quality \(q\) is produced with parental time and market inputs.
    \item Childcare policy changes \(p_x\), availability, reliability, or timing.
    \item The labor response can be participation, hours, earnings, occupation, firm choice, or job quality.
\end{itemize}
\end{frame}

\begin{frame}{Why comparative advantage is not enough}
\begin{itemize}
    \item Wage differences can rationalize some specialization, but they do not explain all gendered care allocation.
    \item Persistent asymmetry may also reflect childcare constraints, bargaining, employer flexibility, and norms.
    \item The empirical task is to avoid labeling every post-birth outcome as preferences or as discrimination.
\end{itemize}
\end{frame}

\begin{frame}{Fertility and child penalties}
\begin{itemize}
    \item Fertility changes care demand and lifecycle timing.
    \item Child penalties trace dynamic post-birth paths for mothers and fathers.
    \item Outcomes include employment, hours, earnings, occupation, firm attachment, promotion, and job quality.
    \item The child penalty is a labor-market object, not a mechanism by itself.
\end{itemize}
\end{frame}

\begin{frame}{Child-penalty event-study object}
\[
Y_{it}^{g} =
\sum_{k \neq -1}\beta_k^{g}\mathbf{1}\{t-B_i=k\}
+\alpha_i+\gamma_t+\varepsilon_{it}
\]
\begin{itemize}
    \item \(B_i\) is first birth and \(k\) is event time.
    \item \(\beta_k^{g}\) traces the dynamic path relative to the pre-birth period.
    \item Best for timing and persistence; not automatically a causal fertility estimate.
\end{itemize}
\end{frame}

\begin{frame}{Fertility effects are a different object}
\[
Y_i = \alpha + \tau K_i + X_i'\delta + u_i,
\qquad K_i \leftarrow Z_i
\]
\begin{itemize}
    \item \(K_i\) is fertility and \(Z_i\) is a fertility shifter such as IVF success.
    \item The target is a causal fertility effect for a specific margin.
    \item External validity depends on which families and fertility margins the shifter moves.
\end{itemize}
\end{frame}

\begin{frame}{Bargaining and control over resources}
\[
\max \lambda U_f(c_f,\ell_f,q)+(1-\lambda)U_m(c_m,\ell_m,q)
\]
\begin{itemize}
    \item Unitary models pool resources and provide a clean benchmark.
    \item Collective models allow resource control and outside options to move \(\lambda\).
    \item A policy effect may be a price effect, income effect, or bargaining effect.
\end{itemize}
\end{frame}

\begin{frame}{Childcare and policy constraints}
\begin{itemize}
    \item Childcare can relax a binding time constraint without changing wages.
    \item Single-parent and coupled households may respond on different margins.
    \item The relevant outcome may be employment, hours, earnings, retention, occupation, firm choice, or job quality.
    \item Policy evidence is strongest when it states the margin moved and the counterfactual.
\end{itemize}
\end{frame}

\begin{frame}{Designs and what they identify}
\scriptsize
\setlength{\tabcolsep}{3pt}
\begin{tabular}{p{0.25\linewidth} p{0.23\linewidth} p{0.18\linewidth} p{0.16\linewidth}}
\toprule
Setting & Object & Method & Main caution \\
\midrule
Birth event study & child penalty path & event time / panel FE & descriptive by itself \\
IVF or infertility & causal fertility effect & IV / quasi-experiment & local margin \\
Childcare reform & constraint relaxation & RCT / DID & multiple channels \\
Resource-control shift & bargaining response & RF / structural & map shift to power \\
Time-use panel & care reallocation & descriptive / panel & weak on earnings dynamics \\
\bottomrule
\end{tabular}
\end{frame}

\begin{frame}{Margin discipline}
\begin{itemize}
    \item Extensive margin: does the worker participate?
    \item Intensive margin: how many hours and which schedule?
    \item Earnings: hours, wages, employment, and career progression combined.
    \item Job quality: flexibility, predictability, commute, promotion, authority, and firm attachment.
    \item Care time: unpaid work can move even when payroll outcomes do not.
\end{itemize}
\end{frame}

\begin{frame}{Norms bridge to Week 5}
\begin{itemize}
    \item Domestic responsibility allocation partly reflects gender norms and traditions.
    \item Norms about caregiving, motherhood, breadwinning, and relative earnings can sustain specialization.
    \item Week 2 flags this mechanism; Week 5 studies norms, beliefs, identity, and institutions directly.
\end{itemize}
\end{frame}

\begin{frame}{Week 2 lab}
\begin{itemize}
    \item Reproduce: synthetic child-penalty event-time summaries.
    \item Diagnose: what the event study shows, and what it cannot identify alone.
    \item Transfer: synthetic IVF-style fertility design and local causal fertility effects.
    \item Optional policy prompt: how childcare availability changes the same household allocation problem.
\end{itemize}
\end{frame}

\begin{frame}{Bridge forward}
\begin{itemize}
    \item Week 3: care responsibilities and fertility timing shape skills, aspirations, and occupational sorting.
    \item Week 5: norms and institutions help explain persistent domestic responsibility allocation.
    \item Later policy weeks ask which interventions move the actual labor margin: care time, employment, hours, earnings, or job quality.
\end{itemize}
\end{frame}

\begin{frame}{Takeaways}
\begin{itemize}
    \item Care work is a production and allocation problem.
    \item Child penalties, fertility effects, bargaining shifts, and childcare policy effects identify different objects.
    \item The labor margins are multidimensional: participation, hours, earnings, occupations, firms, promotions, care time, and job quality.
    \item Week 2 is the household and lifecycle foundation for the rest of the course.
\end{itemize}
\end{frame}

\end{document}
